\chapter{What a Measurement Is}
\label{chap:what-a-measurement-is}

% Closed question: What is the smallest admissible mark a ledger can carry?

\begin{chapteressay}

The nurse wraps a cuff around a patient's arm, closes the valve, and begins to
pump. The cuff tightens. The artery under it is compressed until the blood no
longer passes cleanly with each beat. A dial rises. Then the valve opens by a
controlled amount, the pressure falls, and the nurse listens through a
stethoscope for the first returning sound. A number is read. The cuff continues
to deflate. The sound changes, then disappears. A second number is read.

The patient leaves with a record: 120 over 80.

That record is not blood pressure itself. It is a pair of marks produced by an
instrument under a calibration. The cuff does not touch a metaphysical quantity
called pressure. It applies a counter-pressure to an artery, waits for a
detectable acoustic event, and compares the dial reading at that event against a
scale. The first mark is the pressure at which the artery begins to open during
each pulse. The second mark is the pressure at which the artery remains open
between pulses. The reading is bodily, practical, and familiar, but its logic is
severe: the instrument records two distinguishable events carried by a material
apparatus and read against a reference scale.

That is the structure this chapter names.

A measurement begins when an event becomes admissible for a ledger. Something
happens, but not everything that happens is a measurement. The cuff squeezing
the arm is not yet a measurement. The patient's anxiety is not yet a
measurement. The nurse's hand on the bulb is not yet a measurement. The first
returning Korotkoff sound becomes a measurement only because it crosses the
instrument's threshold, is carried through the stethoscope, is compared against
the gauge, and is written as a mark. The mark has a carrier, a distinction, a
resolution, and a calibration. Remove any one of those and the number loses the
right to be a reading.

The same structure appears in a car's speedometer. The speedometer does not
touch speed. A Hall-effect sensor fires when a magnetic tooth passes. The engine
control unit receives a voltage pulse. The pulse is counted. The count is
scaled by a calibrated wheel circumference and divided by elapsed time. The
dashboard displays a number. Each pulse is like the returning sound in the
blood-pressure cuff: it is not the thing measured, but the admissible event by
which the instrument can carry a claim about the thing measured.

The two examples differ in almost every physical detail. One uses air pressure,
arterial flow, acoustic turbulence, and a human ear. The other uses magnetic
teeth, voltage pulses, a clock, and a dashboard display. Yet both perform the
same grammar. There is an event. There is a carrier. There is a threshold that
decides which events enter the record. There is a calibration that names the
reference before the comparison is trusted. There is a count or ratio of counts.
Only then is there a reading.

This chapter is about the smallest physical structure that can deserve the name
measurement. The later chapters will discuss comparison among instruments,
interpolation between events, classical mechanics, quantum transport, general
relativity, gauge symmetry, and entropy. None of that machinery is available
yet. Before a law can be forced by a record, the record has to exist. Before a
record can exist, a mark must be admissible. Before a mark is admissible, the
instrument must distinguish one event from another and carry the distinction
without confusing it with the world.

The temptation is to begin physics with a number. The thermometer says 98.6, the
speedometer says 60, the ruler says 12 inches, the voltmeter says 5.1 volts. But
the number is late. A reading is the end of a chain, not the beginning. The
physical world presses on the instrument in many ways. The instrument admits
only some of those pressures into its ledger. It does so through a gate: a
threshold, a sensor, a scale, a convention, a calibration. The reading is what
survives that gate.

The closed question of this chapter is therefore deliberately small:

\begin{center}
\emph{What is the smallest admissible mark a ledger can carry?}
\end{center}

The answer cannot be ``a number,'' because a number without a carrier cannot be
read. It cannot be ``an event,'' because an event that cannot be distinguished
from its neighbors cannot be recorded. It cannot be ``a signal,'' because a
signal below threshold is not yet admissible. It cannot be ``a truth about the
world,'' because the measurement may be noisy, approximate, delayed, or
miscalibrated while still being a measurement.

The smallest admissible mark is an event carried as a distinguishable record
under a named calibration. That is the physical atom of this book. The rest of
the volume asks what follows once such atoms are compared, combined,
interpolated, transported, strained, symmetrized, and closed.

\end{chapteressay}

% Phenomenon arcs use: 1/3 setup -> phenomenon box -> 2/3 structural interpretation.

\section{The Event}
\label{se:the-event}

The first obligation of an instrument is not accuracy. It is admission. Out of
all the physical changes occurring around the apparatus, some changes enter the
ledger and others do not. The admitted change is the event.

For the speedometer, the event is the Hall-effect pulse. A ring gear, tone
wheel, or magnetic pickup presents a sequence of distinguishable teeth to a
sensor. When one tooth passes, the local magnetic field changes. The sensor
responds with a voltage pulse. That pulse is the instrument's primitive event.
It is the smallest thing the speedometer can say happened.

The car, of course, is doing vastly more than producing pulses. The tire flexes.
The rubber warms. The road surface deforms. The drivetrain vibrates. The driver
makes tiny corrections. Air moves around the body. None of these happenings is
the speedometer's event. They may matter to other instruments, but they do not
enter this ledger unless they change the pulse record. The speedometer's world
is therefore not the whole physical world. It is the world as filtered by the
event interface.

This is not a defect. It is what makes measurement possible. If every physical
detail entered the record, the record would not be a reading; it would be an
unbounded copy of the world. A finite instrument must have an event boundary. It
must decide what counts. The Hall-effect pulse is precisely such a decision:
tooth passed, event admitted; no tooth passed, no event admitted.

This gives the event two sides. On one side, it is physical: a real voltage
change, generated by a real sensor, caused by a real moving part. On the other
side, it is logical: one new mark in an ordered record. The same pulse belongs
to both registers. If it were not physical, the instrument would not touch the
car. If it were not logical, the instrument would not produce a count.

The nurse's cuff has the same structure. The event is not the entire circulation
of blood through the arm. It is the first audible return of flow at one pressure
and the disappearance of that sound at another. The thermometer has the same
structure. The event is not all molecular motion in the body; it is the
stabilized expansion of a material or the stabilized voltage across a sensor
under a calibration. A Geiger counter has the same structure. The event is not
the radioactive atom's entire history; it is the discharge pulse produced when
ionization crosses the detector's threshold.

In each case, the event is a permitted entry. The instrument does not begin with
a theory. It begins with a gate. The gate is a physical interface whose output
can be counted.

That is why the first mark must be small but not empty. It must be small enough
to be indivisible for the instrument: one tooth, one click, one threshold
crossing, one audible return. But it must not be empty, because an empty mark
would not distinguish the world before from the world after. A mark is the
record of a change the instrument has permission to remember.

The phrase ``permission to remember'' will matter later. In later chapters,
records will be compared, transported, and closed. A comparison can only compare
what has been admitted. A transport can only preserve what the carrier retains.
An entropy argument can only count what the ledger cannot erase. If the event is
wrong, every later law is wrong in the same direction.

For now, the speedometer's first law is humble:

\[
\text{one tooth crossing} \longmapsto \text{one admissible pulse}.
\]

That arrow is not yet speed. It is not yet distance. It is not yet a physical
law. It is the opening of a ledger.

\section{The Carrier}
\label{se:the-carrier}

An event that cannot be carried cannot become a measurement. It may happen, but
it does not persist as a reading.

The speedometer's pulse is carried through a chain. The sensor produces a
voltage change. The wire carries that change to the engine control unit. The ECU
registers the pulse as a digital transition. A counter increments. A clock
records elapsed time. A program scales the count by the calibration constant.
The dashboard display carries the resulting number to the driver.

None of these carriers is the car's speed. The wire is not speed. The counter is
not speed. The dashboard needle is not speed. The display is not speed. The
carrier is the physical support by which the admitted event remains available
for later comparison. Confusing the carrier with the measured quantity is one of
the easiest ways to misunderstand an instrument.

The same count can be carried in many ways. It can be stored as an integer in an
ECU register, displayed as a needle angle, logged as a timestamped data packet,
or printed as a service diagnostic. The physical carrier changes; the carried
event count is meant to remain the same. This is why a measurement record is
never merely a material object. It is a material object under an equivalence
condition. A number on the dashboard and a number in a diagnostic file can be
the same reading if they carry the same admitted events through the same
calibration.

The carrier also imposes its own limits. A frayed wire can corrupt pulses. A
low-resolution display can round the result. A saturated counter can overflow.
A slow display can lag behind a rapid acceleration. These are not failures of
the metaphysical concept of speed. They are failures or limits of the carrier.
The instrument's reading is only as faithful as the carrier's ability to
preserve the admitted distinctions.

This distinction between event and carrier clarifies a common ambiguity. A
speedometer ``reads 60 mph.'' That sentence sounds as if the instrument has
direct access to speed. More literally, the instrument carries a scaled ratio of
event count to time count, and the display presents that ratio in units of miles
per hour. The reading is a carried claim. The claim may be accurate, inaccurate,
delayed, or rounded. It remains a measurement because the carrier makes the
event record available for comparison.

The blood-pressure cuff again helps. The returning sound is carried through
tissue, air, rubber tubing, metal, and the listener's ear. A digital cuff
carries an oscillation pattern through a pressure sensor and a processor. The
event is not the same as the carrier. The carrier is the apparatus that lets the
event survive long enough to be compared with the scale. If the tube leaks, the
event may still occur in the artery, but the instrument fails to carry it
faithfully.

The carrier is therefore the physical answer to a logical problem: how can a
momentary event become part of an ordered record? The answer is not ``by being
true.'' Truth alone does not write. The event must be supported by something
that endures beyond the instant of occurrence. The support may be a mark on
paper, a charge in memory, a mechanical displacement, a human notation, or a
proof term in a formal system. In Volume 2, it is usually a physical apparatus.
In Volume 3, it will often be the compiler's internal record. The grammar is the
same.

This is the smallest form of abstraction in the physical gauge. Abstraction
does not mean leaving matter behind. It means recognizing that the same event
record can be carried by different material systems while preserving the
distinction the measurement needs. The carrier makes portability possible.

If the event is the admitted change, the carrier is the instrument's memory of
that change.

\section{Distinguishability}
\label{se:distinguishability}

A carrier is useful only if it can keep events apart. The third obligation of a
measurement is distinguishability.

For the speedometer, distinguishability begins with the teeth. Suppose the tone
wheel has thirty-six magnetic teeth per revolution. At one wheel revolution,
the sensor can admit thirty-six pulses. If the wheel turns once and the
instrument registers thirty-five pulses, one event has been lost. If it
registers thirty-seven, one event has been invented or duplicated. The count is
trustworthy only when adjacent teeth generate adjacent but separate events.

At 60 mph, a typical passenger-car wheel might rotate on the order of thirteen
times per second, depending on tire circumference. A thirty-six-tooth sensor
then sees roughly 470 pulses per second. The exact number is not important here.
What matters is the spacing. The sensor must recover from one pulse before the
next tooth arrives. The wire and ECU must distinguish one voltage transition
from the next. The clock must distinguish one time interval from another. If the
apparatus smears two adjacent pulses into one, the ledger loses a mark.

Distinguishability is not the same as difference in the world. Two physical
states can be different without being distinguishable by a given instrument. A
car traveling at 60.000 mph and a car traveling at 60.001 mph are different
states if the units are taken literally. But the dashboard may not distinguish
them. If the display rounds to the nearest mile per hour, both states produce
the same visible reading. They are different for a better instrument, perhaps,
but not for this ledger.

This is the beginning of physical humility. The instrument does not say what
differences exist in the world. It says which differences it can carry. A
measurement record therefore has a grain. Below that grain, differences may
exist, but they are not events for this instrument. Above that grain, they can
enter the record.

This is why the physical gauge refuses to begin with a continuum. The continuum
may appear later as a smooth limit, but no finite instrument begins there. A
finite instrument begins with distinguishable alternatives. Pulse or no pulse.
Click or no click. Reading changed or reading unchanged. The continuum is a
mathematical completion of a record, not the primitive form of the record.

Distinguishability also explains why a record can be useful without being
omniscient. A speedometer does not record every microscopic change in tire
motion. It records enough distinguishable tooth crossings to produce a stable
speed estimate. A thermometer does not record every molecular collision. It
records enough aggregate sensor response to produce a temperature reading. A
voltmeter does not record every electron. It records a distinguishable potential
difference within its resolution.

The formal shape is simple. For an instrument, two admitted events must satisfy
an operational distinction:

\[
e_i \neq e_j \quad \text{as records, whenever the instrument can separate them.}
\]

This is not a metaphysical inequality. It is a claim about the ledger. The
instrument can write one mark for $e_i$ and another mark for $e_j$, and those
marks do not collapse into the same entry.

Distinguishability is therefore the first physical source of order. Once events
can be kept apart, they can be counted. Once they can be counted, they can be
compared. Once they can be compared, ratios can be formed. A speedometer's
reading begins not with velocity, but with the ability to say: this pulse is not
that pulse.

That is the smallest ordering a physical instrument needs.

The following phenomenon is placed here because it clarifies the cost of
moving from repeated events to physical expectation. Distinguishable marks do
not by themselves produce certainty about unmeasured cases. They produce a
record, and the record always carries induction pressure.

\begin{phenom}{Hume Effect}
\label{ph:hume-effect}

\PhStatement A finite measurement record cannot contain the law that extends it
to all future cases. It can contain repeated marks, stable comparisons, and a
calibration history, but the inference from those marks to a universal rule is
an additional obligation.

\PhOrigin David Hume's problem of induction: no finite sequence of observed
regularities logically entails that the same regularity will continue in the
next unobserved case.

\PhObservation A speedometer can be tested repeatedly against a calibrated track
and found to agree each time. Those trials justify trust in the instrument, but
the next reading is still an inference from the record and the calibration
procedure, not a logical consequence contained inside any single pulse.

\PhConstraint The instrument may only use admitted events and named calibration
standards. It may not smuggle in an unrecorded guarantee that future cases will
match past cases.

\PhConsequence Repeatability is necessary for measurement, but repeatability is
not omniscience. A physical law must be compatible with the finite record and
its calibration history; it is not literally stored in the first mark.

\end{phenom}

The Hume Effect does not make measurement impossible. It prevents a false
inflation of measurement. A finite ledger can support inference, but only after
its limits are named. The speedometer's pulse record supports a stable speed
reading because the sensor, counter, clock, and calibration have been tested and
because future readings are constrained by the same apparatus. But the record
does not contain all possible future roads. It contains marks under a rule.

This will matter later when laws appear. The book will argue that physical laws
are forced by admissible measurement records. That does not mean one pulse
contains Newtonian mechanics, general relativity, or thermodynamics. It means
that once a record is required to be finite, distinguishable, repeatable,
calibrated, transportable, and closed, many candidate descriptions become
inadmissible. The law is forced by the full discipline of the ledger, not by a
single isolated mark.

\section{Threshold and Noise}
\label{se:threshold-and-noise}

Every instrument has a threshold. Below it, the world may change, but the ledger
does not.

For the speedometer, the low-speed threshold is familiar. A car can creep
forward in a driveway while the dashboard still reads zero. The wheels move.
The drivetrain turns. The sensor may even see occasional pulses. But the update
rate, display rounding, and filtering logic may decide that no stable nonzero
speed should be shown. The instrument is not denying motion. It is saying that
the motion has not entered the displayed record as a distinguishable speed.

Thresholds are often mistaken for errors. Sometimes they are. A broken sensor
can fail to register an event that should have been admitted. But a threshold as
such is not an error. It is the boundary that makes the instrument finite. A
sensor without threshold would admit every fluctuation as an event. The record
would drown in noise.

Noise is not merely an annoyance added after the ideal measurement. Noise is the
reason thresholds are needed. The sensor voltage fluctuates. The wheel vibrates.
The tire slips. The road surface produces small irregularities. The electrical
system carries interference. The ECU must decide which changes are pulses and
which changes are background. A measurement exists only after that decision.

This gives threshold and noise a paired structure. If the threshold is too low,
noise becomes false event. If the threshold is too high, real events disappear.
The instrument's usable range lies between these failures. A good speedometer
does not eliminate noise from the universe. It chooses a gate through which the
event of interest can pass more reliably than the background fluctuations.

The same pattern appears in static friction. A heavy box rests on the floor. A
small force is applied. Nothing moves. A slightly larger force is applied. Still
nothing moves. At some threshold, the box slips. The applied force did not
become real only at the moment of slipping; it was real before. But before the
threshold, it did not enter the motion ledger as displacement. The table and box
absorbed it as strain. The event ``box moved'' was not admitted.

That structure is important enough to name.

\begin{phenom}{da Vinci-Coulomb Effect}
\label{ph:davinci-coulomb-effect}

\PhStatement Static friction creates a physical threshold: applied burden can
increase while the displacement record remains unchanged, until a critical
condition is crossed and motion becomes admissible.

\PhOrigin Leonardo da Vinci described empirical laws of friction in notebooks;
Charles-Augustin de Coulomb later developed systematic friction laws relating
frictional resistance to normal load and material contact.

\PhObservation A box pushed gently across a table does not move at first. The
applied force is real, but the ledger of position carries no new mark. Once the
static-friction threshold is exceeded, the box slips and displacement becomes a
recorded event.

\PhConstraint The measurement record may not count sub-threshold burden as
motion. It can record strain, applied force, or deformation if equipped to do
so, but the displacement ledger admits a mark only when distinguishable motion
occurs.

\PhConsequence Physical measurement is gated. The absence of a mark does not
mean the absence of all physical change; it means the relevant event has not
crossed the instrument's admissibility threshold.

\end{phenom}

The speedometer's threshold is less dramatic than the slipping box, but the
logic is the same. Below the pulse rate that the display treats as stable, the
car's motion may be physically real and practically visible, but it is not a
speedometer reading. The ledger remains at zero. That zero is not a metaphysical
claim of perfect rest. It is an instrument claim: no admissible speed event has
crossed the display threshold.

This distinction saves the later argument from a common mistake. If the
instrument reads zero, one might say ``nothing happened.'' But that is too
strong. The correct statement is: nothing happened that this instrument, under
this calibration and threshold rule, admitted as a mark in this ledger. That is
weaker, but it is honest. Physical law must be built from honest claims, not
inflated ones.

Thresholds also explain why improving an instrument changes the record without
changing the past. A higher-tooth-count sensor, a faster update rate, or better
noise filtering may admit events that the old speedometer missed. The old
record was not false simply because the new record is finer. It was a coarser
ledger. Refinement adds distinctions; it does not magically insert them into
the earlier instrument.

This is the physical beginning of scale. A measurement has a scale because the
instrument has a threshold. Below the scale, differences are not carried by the
record. Above the scale, differences can become events. Every later continuum
description must remember that it is a smooth account of records that began in
thresholded marks.

\section{Calibration Before Measurement}
\label{se:calibration-before-measurement}

The event, carrier, distinction, and threshold still do not make a trusted
measurement. One more structure is needed: calibration.

The speedometer counts pulses. To turn those pulses into speed, it must know how
far the car travels per pulse. That requires the wheel circumference, gear
ratio, and tooth count. If the tire radius changes, the calibration changes. A
larger tire travels farther per revolution than the original tire. The same
pulse count per second now corresponds to a higher road speed. If the
speedometer still uses the old circumference, the display reads too low.

This is not a problem that can be solved by counting more carefully. The pulse
count may be perfect. The clock may be perfect. The carrier may be perfect. If
the reference circumference is wrong, the reading is wrong. Calibration is not a
decorative correction after measurement. It is the condition under which the
count becomes a physical quantity.

The simplest speedometer computation has the form

\[
\text{speed}
  =
  \frac{\text{pulses}}{\text{time}}
  \cdot
  \frac{\text{wheel circumference}}{\text{pulses per revolution}}.
\]

The first factor is a record: pulses per unit time. The second factor is a
calibration: distance per pulse. The reading is the product. Without the second
factor, the instrument has a frequency, not a speed.

Calibration must be performed before the reading is trusted. This is a
practical rule in laboratories and garages, but it is also a logical rule. A
reading without a named reference cannot say what it is a reading of. The
reference may be a tire circumference, a gauge block, a cesium transition, a
standard resistor, a kilogram realization, or a calibration certificate from a
metrology lab. Whatever its material form, it plays the same role: it names the
scale against which the admitted marks are compared.

This is why the Lean calibration module matters even in the physical volume. In
the code, the calibration hides its raw count and exposes only admissible
questions through an interface. The physical analog is straightforward. The
driver need not inspect the entire calibration history of the speedometer while
driving. The driver sees a public reading. But the reading is public only
because a private calibration standard has already been carried into the
instrument.

Calibration also prevents a subtle confusion between repeatability and truth. An
uncalibrated instrument can be perfectly repeatable. A speedometer with the
wrong tire radius may read 55 every time the car is actually traveling 60. It is
stable, but stably wrong. Repeatability shows that the instrument preserves its
own comparison. Calibration names the external reference that lets the
comparison mean what it claims.

The same pattern appears in the blood-pressure cuff. The pressure gauge must be
calibrated against a reference. The cuff size must fit the arm. If the cuff is
too small, the pressure reading may be too high; if too large, too low. The
returning sound can be heard correctly and the dial read carefully, while the
measurement remains biased by a bad reference. Calibration is therefore not an
optional act of precision. It is the instrument's right to speak in units.

The smallest admissible mark of this chapter is now complete. It is not merely a
physical change. It is not merely a signal. It is not merely a counted pulse. It
is an admitted event, carried by an apparatus, distinguished from neighboring
events, gated against noise, and read under a named calibration.

Only such a mark can enter the ledger that the rest of the book will use.

\begin{bridge}

The chapter's question was: what is the smallest admissible mark a ledger can
carry?

The answer is now precise enough to move. A mark is an admitted physical event
carried as a distinguishable record under a threshold and a calibration. The
speedometer's Hall-effect pulse is such a mark only because the sensor admits
it, the wire and ECU carry it, the tooth spacing distinguishes it from its
neighbors, the display logic gates it against noise, and the tire circumference
calibrates it as distance per pulse.

This is still not a physical law. It is the material needed before physical law
can begin. A single mark can enter a ledger, but physics requires comparison:
one mark against another, one instrument against another, one reading against a
reference, one repeated trial against the next. Chapter 2 therefore asks how an
instrument can compare records and repeat the result without pretending that
the comparison is the thing measured.

\end{bridge}

\begin{coda}{The Turnstile at a Transit Station}

A commuter approaches a transit gate during the morning rush. The station is
loud, crowded, and full of motion. People shift bags from one shoulder to
another. Shoes scrape. Phones light up. Coins drop. A train arrives below the
platform and the floor trembles. The turnstile ignores almost all of it.

The gate waits for one event: a valid fare tap.

The commuter brings a card near the reader. At first nothing happens. The card
is still too far away, or the angle is wrong, or the signal is too weak. The
passenger is physically present, the card exists, the intention to enter is
clear, but the instrument has admitted no event. Then the card crosses the
reader's threshold. The chip responds. The reader receives the credential. The
fare system compares it against the current rule. The light turns green. The
bar unlocks. The passenger pushes through.

One passage has entered the ledger.

The turnstile does not know the passenger in any rich sense. It does not know
where the passenger is going, whether the trip is urgent, whether the passenger
is late, whether the weather outside is cold, or whether the train below is
already full. It knows the event it is built to admit. A credential crossed the
reader's threshold; the fare was valid under the current calibration; the gate
opened; one body passed through the arms. That is the mark.

The carrier is distributed. The card carries an account identifier. The reader
carries the tap to the station controller. The controller carries the decision
to the gate motor. The turnstile arms carry the decision into physical space by
locking or unlocking. The central fare database carries the transaction into
the day's record. No single material object is the passenger's passage. The
passage is carried through the system by a chain of supports.

Distinguishability is just as important. The gate must distinguish one valid tap
from another. If two commuters press forward at once, the arms admit only one
body per cycle. If the reader sees two cards at the same time, it may reject the
tap rather than invent a fare event. If a passenger taps twice too quickly, the
system may suppress the second tap as a duplicate. The turnstile is not trying
to describe the entire crowd. It is trying to maintain a ledger in which one
admissible passage is not confused with another.

Threshold and noise are everywhere. A card in a pocket near the reader may not
count. A phone case may block the signal. A weak credential may produce an
ambiguous response. The station's electrical systems may add interference. The
gate is built to ignore sub-threshold noise because otherwise the morning rush
would become a storm of false passages. The absence of a green light does not
mean no human being is present. It means no admissible fare event has crossed
the gate's threshold.

Calibration enters through the fare rule. The same card may be valid at one
time and invalid at another. The fare may be different for a student, a senior,
a transfer, a rush-hour trip, or a monthly pass. Before the gate can decide, the
current fare table must be named. Without that reference, the tap is only a
signal. With the reference, it becomes a fare event. The calibration certificate
is not a physical ruler, but it performs the same function as the speedometer's
tire circumference: it lets the raw event be interpreted as a reading.

The turnstile therefore performs the chapter's machinery without using any of
the chapter's physical examples. The event is the valid tap. The carrier is the
card-reader-controller-gate-database chain. Distinguishability is one passage
per admitted cycle. Threshold is the reader's gate between noise and signal.
Calibration is the fare rule. The ledger is the day's entry count.

If the passenger slips behind someone else without tapping, the body passes but
the ledger has no admissible mark. If the reader charges the card but the gate
does not open, the ledger and the physical passage disagree. If the fare table
is wrong, the gate can be perfectly repeatable and still wrong. If the database
loses the transaction, the event happened but was not carried. Each failure is a
different failure because each part of the measurement structure is different.

The turnstile is not a speedometer. It measures no velocity, contains no tire
radius, and knows nothing of Hall-effect sensors. Yet it proves the portability
of the chapter's answer. The smallest admissible mark is not a number floating
free of matter. It is an event admitted by an instrument, carried by a material
support, distinguished from its neighbors, gated against noise, and interpreted
under a calibration.

At the end of the day, the station manager may ask how many passengers entered.
The answer is not the number of people who wanted to enter, nor the number of
feet that crossed the station floor, nor the number of bodies visible on a
security camera. It is the number of admissible passages the turnstile carried.

That is a measurement.

\end{coda}
